\documentclass[11pt,titlepage]{article}

\usepackage[french,american]{babel}
\usepackage[utf8]{inputenc}
\usepackage[T1]{fontenc}
\usepackage{lmodern}
\usepackage{amsmath,amsfonts,amssymb}
\usepackage{graphicx}
\usepackage[a4paper]{geometry}           % See geometry.pdf to learn the layout
%\geometry{landscape}                % Activate for for rotated page geometry
%\usepackage[parfill]{parskip}    % Activate to begin paragraphs with an empty line rather than an indent
\usepackage{graphicx}
\usepackage{amssymb}
\usepackage{epstopdf}
\usepackage{color}
\usepackage[tt]{titlepic}
\usepackage{fancyhdr}
\usepackage{secdot}
\usepackage{booktabs}
\usepackage{rotating}
\usepackage{multirow}
\usepackage[style=ieee, backend=bibtex]{biblatex}

\topmargin=-0.45in      %
%\evensidemargin=0in     %
\oddsidemargin=-0.1in      %
\textwidth=6.8in        %
\textheight=9.0in       %
%\headsep=0.25in         %

%======================= Titlepage =========================================================
\DeclareGraphicsRule{.tif}{png}{.png}{`convert #1 `dirname #1`/`basename #1 .tif`.png}
\titlepic{\includegraphics[width=17cm]{figures/research_plan_pic.pdf}}

\title{\textbf{Annual progress report \textcolor{blue}{2014}}} %Change Year

\author{\large{\textbf{Damar Wicaksono} }\\\\ 
\large{\textbf{Prof. Andreas Pautz}} \\\
\large{\textbf{Omar Zerkak}} \\\\\\\
\Large{\textbf{Bayesian Uncertainty Quantification}} \\\\
\Large\textbf{of Physical Models} \\\\
\Large\textbf{in Thermal-Hydraulics System Code}}

\date{30.11.2014}				%activate to remove date
%========================================================================================

\bibliography{report_2014.bib}

\begin{document}
%========================  Top header ======================================================
\pagestyle{fancy} \pagenumbering{arabic} \setcounter{page}{1}
\addtolength{\headheight}{\baselineskip}
\newcommand{\ffont}{\fontsize{8}{8}\selectfont}
\lhead{ECOLE DOCTORALE\\ \textit{DOCTORAL SCHOOL}}
\chead{PROGRAMME DOCTORAL EN PHYSIQUE\\ \textit{DOCTORAL PROGRAM IN PHYSICS}}
\rhead{\bfseries\ffont  \includegraphics[width=55pt]{figures/EPFL_LOG.pdf}}
\renewcommand{\headrulewidth}{0.4pt}
\maketitle
%=========================================================================================
\textcolor{white}.
\section{Objectives of research}

The objective of the research is to quantify the uncertainty of 
physical model parameters implemented in a thermal-hydraulics system code.
The physical models concerned are the ones describing 
the phasic interaction (heat and momentum exchange) in a complex 
multiphase flow during reactor transient especially for the
\textit{reflood} phase after a loss-of-coolant-accident. 
These models parameterized either by a physical or tuning 
parameters.

Conforming with the practice of statistical uncertainty propagation widely 
adopted in the field of nuclear engineering, probability theory is used to 
quantify the uncertainties related to these parameters with the results in 
a form of density function and/or its approximation. 
The derivation of this density is posed as an inverse statistical problem 
following a bayesian approach as the parameters themselves are not directly 
observable.
To this end, a methodology will be developed based on probabilistic 
modeling combined with the available relevant experimental data taken 
from various separate effect test facilities.

\section{Work achieved in the past year (state of research)}

\section{Current state of work}

The current state of research along with the achievements during the year are 
summarized in the table below.
                   analysis and
                               novel data analysis methodology (NUTHOS-14)
                          1 abstract for the NURETH-16 submitted 
                                             
                     
\begin{table}[!h]
    \begin{center}
        \begin{tabular}{  c  p{4cm}  p{2.5cm} p{8cm}}
        \toprule
        \bf{Phase} & \bf{Task Description} & \multicolumn{2}{c}{\bf{Outcome}}  \\
                   &                       & \multicolumn{1}{c}{\bf{Planned}} & \multicolumn{1}{c}{\bf{Current}}       \\\toprule
% Phase 1         
                     \multirow{5}{*}{1}
                       & Comprehensive reviews of post-CHF flow  & 
                         \multirow{3}{*}{$\bullet$ 1 Tech. Rep.} 
                          & $\bullet$ Relevant parameters externalized \\ 
                       & closure laws in TRACE and externalization & 
                           & $\bullet$ PSI contribution to PREMIUM benchmark  \\
                       & of important model parameters & & \\              
                   \midrule     

% Phase 2                     
                     2
                   & Global sensitivity analysis of important & & \\
                   & reflood model parameters & & \\               
                   &  &  $\bullet$ 1 Technical Report & \\
                   &  &  $\bullet$ 1 Conference Paper & \\
                     \\\midrule
% Phase 3         
                     3
                   & Calibration of TRACE reflood model & & \\
                   & based on FEBA test facility & & \\               
                   & & $\bullet$ 1 Technical Report  & \\
                   & & $\bullet$ 1 Journal Article   & \\               \\\midrule
% Phase 4
                     4
                   & Calibration of TRACE reflood model & & \\
                   & based on another test
                     facility (To be Decided)               
                   & 1 Conference Paper                            
                    \\\midrule
% Phase 5
                     5
                   & Consolidation of the calibration results & & \\
                   & based on 2 facilities and validation & & \\
                   & based on another reflood test facility & & \\         
                   & 1 Journal Article            
                    \\\midrule 
% Phase 6
                     Phase 6
                   & Thesis write-up               
                   & Thesis                
                   & $-$\\
                     \bottomrule    
    \end{tabular}
\end{center}
\end{table}


\section{Calendar of upcoming work}


\pagebreak

\section{Scientific publications, conference contributions, and presentations}

\nocite{Wicaksono2014a}
\nocite{Wicaksono2014b}
\nocite{Wicaksono2014c}
\nocite{Wicaksono2014d}
\nocite{Wicaksono2014e}
\printbibliography[heading=none]

\section{Other activities and remarks}

\subsection{Nuclear Computation Lab (ETH-531)}

A teaching activity during the year 2014 was carried out for Nuclear Computation Lab 
course given at PSI.
The course was part of the block courses offered to the EPFL/ETHZ Nuclear 
Engineering Master Program
Single (full) day was dedicated for lecture session followed by computer 
simulation session for the t
12 students participated in the course with 

Prior to the course, the available lecture slides was fully revised for a better structure and with some examples.
 

\end{document}
\newpage
\textcolor{white}.\\\\
\noindent \textbf{Comments by the thesis advisor(s)}\\\\
...................................................................................................................................................................\\\\
...................................................................................................................................................................\\\\
...................................................................................................................................................................\\\\
...................................................................................................................................................................\\\\
...................................................................................................................................................................\\\\
...................................................................................................................................................................\\\\
...................................................................................................................................................................\\\\
...................................................................................................................................................................\\\\
...................................................................................................................................................................\\\\


\section*{\underline{Signatures}\\}
\noindent \textbf{Thesis advisor}\hspace{6.25cm}..............................................................................\vspace{0.5cm}\\

\noindent \textbf{Thesis co-advisor}\hspace{5.7cm}..............................................................................\\
\noindent  (if officially nominated)\\\\

\noindent \textbf{With his/her signature, the candidate confirms that he/she took note of the above comments.}\\\\
\textbf{Candidate}\hspace{7cm}..............................................................................\vspace{0.5cm}\\

\noindent \textbf{Date}\hspace{8.05cm}..............................................................................\\

\vspace{0.8cm}
\begin{center}
\textbf{Up to 5 pages in total}
\end{center}
\begin{center}
\textbf{To be returned to:} \textcolor{green}{EPFL-EDPY / Station 3 / 1015 Lausanne / Suisse}
\end{center}
\end{document}  